% =====================================================================
% Capítulo 1. Introducción
%
% Esqueleto de trabajo. Cada comentario indica qué va en cada sitio y de
% qué ficha de notas/ sale el material. Borrar los comentarios según se
% vaya redactando.
%
% Extensión orientativa: 10-12 páginas contando figuras.
% La guía marca Introducción como sección prioritaria junto al Resumen y
% las Conclusiones: es lo que el tribunal lee seguro.
% =====================================================================

\chapter{Introducción}
\label{cap:introduccion}

Los escáneres intraorales 3D son dispositivos que permiten capturar la geometría y el color de la arcada dental de un paciente, generando una malla tridimensional detallada. 
Estos escáneres se utilizan para planificar tratamientos de ortodoncia, prótesis y otras intervenciones en odontología. Su uso se ha vuelto común, y cada vez más clínicas 
dentales realizan escaneos intraorales 3D de sus pacientes \cite{pendiente}.

Además de su aplicación en el ámbito clínico, estos escaneos tienen un gran potencial en la identificación forense, ya que permiten comparar registros \textit{antemortem} 
y \textit{postmortem}, lo cual puede ser fundamental para establecer la identidad de personas fallecidas. Cuando existen buenos registros previos, la odontología forense 
resulta clave en la identificación de víctimas, ya que llega a identificar en torno al 60\,\% de ellas, y contribuye en otro 30\,\% de los casos \cite{forrestForensicOdontologyDVI2019}.
La literatura forense apunta a que se podrán comparar modelos 3D de manera automatizada \cite{forrestForensicOdontologyDVI2019}. Esto ayudaría en grandes catástrofes o 
accidentes donde se requiera identificar a muchas personas. Además, como cada vez más clínicas escanean a sus pacientes, crece el número de registros \textit{antemortem} disponibles, 
que podrían utilizarse en el futuro para la identificación forense.

Por eso, este trabajo se enfoca en segmentar y etiquetar cada diente, considerando ausencias, apiñamiento y brackets, manteniendo la coherencia de la numeración sin asumir que la arcada
está completa, con vistas a su uso en contextos forenses y clínicos.

\section{Descripción del problema}
\label{sec:descripcion-problema}

Tradicionalmente se tomaba una impresión física de la arcada dental del paciente con un material blando y luego se creaba un modelo de escayola. Las técnicas de impresión convencionales 
presentan problemas de incomodidad para el paciente, deformación del material y almacenamiento físico de los modelos \cite{leeClinicalApplicationIntraoral2024}. El escáner intraoral 
suprime la necesidad de realizar la impresión, lo que incrementa la comodidad del paciente y crea un modelo digital que puede ser guardado y revisado de manera indefinida.
Respecto a su exactitud, al evaluar los escáneres intraorales en relación con la impresión de alginato en individuos con dentición mixta, es decir, niños que tienen tanto dientes de leche 
como dientes permanentes, las variaciones promediaron 0,022 mm, lo cual se considera dentro de un rango clínico aceptable \cite{leeClinicalApplicationIntraoral2024}.

De esta manera, lo que solía ser un objeto tangible, que podía extraviarse o deteriorarse con el paso del tiempo, ahora se ha transformado en un archivo que puede guardarse en formato
digital y revisarse en cualquier momento, convirtiéndose en una herramienta para la identificación en investigaciones forenses, dado que puede confrontarse con registros
\textit{postmortem} futuros.

Además, se puede hacer el escaneo sin necesidad de retirar los brackets, permitiendo así registrar la dentadura del paciente mientras se lleva a cabo el tratamiento ortodóntico
\cite{leeClinicalApplicationIntraoral2024}. Sin embargo, los brackets quedan adheridos a la superficie del diente como piezas que no forman parte de él, ocultan su verdadera forma y 
complican su segmentación y etiquetado \cite{ben-hamadou3DTeethSeg223DTeeth2023}.

El tipo de imagen más utilizado en el campo dental es la ortopantomografía (OPT), una radiografía panorámica que captura en una sola imagen bidimensional ambas arcadas
y sus estructuras adyacentes \cite{barcoDeteccionSegmentacion2025}. Al ser una proyección plana, las estructuras se superponen y, además, la imagen presenta distorsiones, ya que la 
ampliación no se distribuye de manera uniforme a lo largo de toda la imagen \cite{barcoDeteccionSegmentacion2025}. Por otro lado, el escaneo intraoral proporciona información 
complementaria: aunque no revela estructuras internas como las raíces o el hueso, representa la superficie real de cada diente en tres dimensiones, a escala y con color, sin superposiciones
y evitando la exposición del paciente a radiación \cite{forrestForensicOdontologyDVI2019}. Investigaciones anteriores del mismo grupo se centraron en la detección, segmentación e 
identificación de dientes en imágenes OPT mediante aprendizaje profundo \cite{barcoDeteccionSegmentacion2025}. Se observó que los errores se concentraban en los molares, dado que son los 
dientes que más frecuentemente están ausentes, lo que provoca que el modelo tenga menos referencias anatómicas para diferenciar estos dientes \cite{barcoDeteccionSegmentacion2025}. 
El presente trabajo lleva el desafío al ámbito tridimensional y se enfoca precisamente en esta complejidad: asignar a cada diente su número FDI de manera coherente con el resto de la 
arcada, aun en casos de dientes ausentes.

Aparte de los brackets, el escaneo intraoral enfrenta otras dificultades propias. En la malla, no hay un límite geométrico bien definido que separe el diente de la encía, ya que
forman una superficie continua, por lo quedelimitar cada diente sea complicado \cite{cuiTSegNetEfficientAccurate2021}. Cuando los dientes están apiñados, las divisiones entre
dientes adyacentes no son evidentes, y así, dos dientes cercanos pueden ser confundidos como uno solo \cite{ben-hamadou3DTeethSeg223DTeeth2023}.

% ACTUALIZAR con las 6 partes: 4,7 %, 35 % e histograma (src/conteo_dientes.py)
La falta de dientes presenta un problema diferente, relacionado con la numeración. En el sistema de notación FDI, cada número está vinculado a una ubicación específica en la arcada, por
lo que si un diente está ausente, su número no se cambia: si el segundo premolar (35) no está, el diente adyacente al primer premolar continúa siendo el primer molar (36). Si se 
utilizara un método que numerara los dientes según su secuencia de aparición, el primer molar recibiría el número 35, y a partir de ahí se desplazaría la numeración de todos los dientes
siguientes. Esta situación no es inusual. En los 300 escaneos revisados de Teeth3DS, únicamente un 4{,}7\,\% muestra la arcada completa de 16 dientes, y el 35\,\% presenta al menos una 
ausencia aun sin contar los terceros molares (Figura~\ref{fig:histograma-dientes}).

\begin{figure}[htbp]
    \centering
    \includegraphics[width=0.75\textwidth]{imagenes/histograma_dientes.png}
    \caption{Número de dientes presentes por escaneo en los 300 escaneos de la primera de las seis partes de Teeth3DS. Solo el 4{,}7\,\% contiene la arcada completa de 16 dientes.}
    \label{fig:histograma-dientes}
\end{figure}

\subsection{Identificación forense y odontología clínica}
\label{subsec:forense-clinica}

La identificación de víctimas de catástrofes (DVI, en inglés) es el proceso por el cual se confirma quiénes son los individuos que han perdido la vida en 
un evento con gran cantidad de víctimas, especialmente cuando su estado impide el reconocimiento a simple vista. En estas situaciones, los dientes se convierten en una herramienta 
fundamental: resisten condiciones extremas y su forma, disposición y los tratamientos que han recibido son característicos de cada persona \cite{forrestForensicOdontologyDVI2019}.

El método se basa en la comparación de dos perfiles dentales. El perfil \textit{antemortem} se elabora utilizando la información de personas desaparecidas, que conservan sus dentistas: 
fichas escritas, odontogramas, radiografías y, cada vez más, tomografías y escaneos en 3D. Las imágenes se consideran más fiables que los registros escritos, ya que 
provienen directamente del paciente \cite{forrestForensicOdontologyDVI2019}. El perfil \textit{postmortem} se genera mediante la evaluación de cada víctima de acuerdo con las directrices 
de INTERPOL, que requieren que se registre la condición de cada diente, ya sea que esté presente o ausente, además de radiografías y fotografías. Cuando hay un gran número de víctimas,
la comparación se apoya en programas como DVISys, que ordena a los posibles candidatos para su evaluación posterior por un especialista \cite{forrestForensicOdontologyDVI2019}.

Para que esta comparación sea efectiva, es fundamental que ambos perfiles se refieran a cada diente de forma uniforme. Hay múltiples sistemas de numeración dental y no todos los países
emplean el mismo, por lo que INTERPOL sugiere el uso de la notación de la Federación Dental Internacional (FDI), la cual otorga a cada diente un número compuesto por dos dígitos
\cite{forrestForensicOdontologyDVI2019}. Un diente que esté mal identificado en cualquiera de los perfiles puede crear una inconsistencia que podría excluir a la persona correcta o
favorecer a una incorrecta. Esto resalta la necesidad de que un sistema automático no solo separe los dientes, sino que también les asigne el número adecuado.

\subsection{Segmentación y etiquetado de los dientes}
\label{subsec:segmentacion-etiquetado}

% La tarea concreta: separar cada diente de la encía (segmentación) y
% asignarle su número FDI (etiquetado). Distinguir bien las dos cosas,
% porque buena parte de la literatura solo hace la primera (ficha 04).
%
% Introducir aquí la notación ISO-3950/FDI solo lo justo; el desarrollo
% va en Fundamentos teóricos.
%
% Figura 1.2: diagrama general del sistema (entrada -> segmentación ->
% etiquetado -> salida). Equivalente a la Figura 1.2 de Barco.

\subsection{Solución propuesta}
\label{subsec:solucion-propuesta}

% Qué se propone, en positivo (la guía insiste en no abrir con lo que el
% método no hace).
%
% El hueco, tal como quedó tras leer TSegNet y DArch (ficha 05, punto 6):
% etiquetar con coherencia global sin asumir que la arcada está completa.
% - Los métodos que usan el contexto de la arcada asumen dentición
%   completa o reordenan secuencias sin huecos.
% - Los datos lo desmienten: 13.3 dientes por escaneo de media en
%   Teeth3DS frente a los 16 de una arcada completa (ficha 03).
%   IMPORTANTE: sustituir por la cifra que salga del script propio de
%   conteo en cuanto se tenga acceso a los datos.
% - La propuesta: asignación conjunta de los números FDI (método
%   húngaro), que no repite números y admite posiciones vacías.

\section{Motivación}
\label{sec:motivacion}

% ¿Por qué es relevante? Dos bloques.
%
% 1. Casos concretos donde la identificación forense hace falta. Barco
%    usa una lista de catástrofes con cifras y referencias; sirve de
%    modelo, pero conviene buscar las fuentes propias y no copiar su
%    selección.
% 2. Por qué en 3D y por qué ahora: el escáner intraoral se ha
%    generalizado en la clínica, así que el volumen de datos antemortem
%    3D crece cada año (ficha 02).
%
% La guía recomienda un histograma de publicaciones del campo (Scopus o
% similar) indicando la consulta empleada y la fecha de la búsqueda.
% Figura 1.3.

\section{Objetivos}
\label{sec:objetivos}

% Un objetivo principal y varios parciales. En las Conclusiones hay que
% volver sobre ellos y decir en qué grado se han alcanzado, así que
% conviene que sean comprobables.
%
% Principal (borrador a pulir con los tutores):
%   Desarrollar un sistema de segmentación y etiquetado de dientes en
%   escaneos intraorales 3D que mantenga la coherencia de la numeración
%   sin asumir que la arcada está completa.
%
% Parciales, uno por línea:
%   - Estudiar el estado del arte en segmentación y etiquetado dental 3D.
%   - Caracterizar Teeth3DS, midiendo cuántos dientes faltan y con qué
%     frecuencia.
%   - Implementar la línea base y la propuesta de asignación conjunta.
%   - Evaluar con métricas que midan la coherencia de la numeración, no
%     solo el acierto por diente (ficha 05, punto 6.3).

\section{Planificación del proyecto}
\label{sec:planificacion}

\subsection{Planificación temporal}
\label{subsec:planificacion-temporal}

% 12 ECTS x 25 h = 300 horas. Repartirlas por etapas y dibujar un
% diagrama de Gantt sencillo (Figura 1.4). Etapas previsibles:
% revisión bibliográfica, acceso y exploración de datos, implementación,
% experimentación, y escritura en paralelo durante todo el proyecto.
%
% Dejar esta subsección viva: al final del TFG hay que añadir la
% planificación real y justificar las desviaciones. La rúbrica puntúa
% explícitamente identificar y justificar desviaciones.

\subsection{Planificación económica}
\label{subsec:planificacion-economica}

% Horas de trabajo a precio de junior, más amortización del material:
% portátil, GPU o servidor de la universidad, licencias. Se presupuesta
% la parte amortizada, no el coste total del equipo.

\section{Uso de herramientas de inteligencia artificial}
\label{sec:uso-ia}

% La guía de Mesejo pide indicar con claridad qué herramientas de IA
% generativa se han usado, de qué modo y en qué partes del trabajo.
% Redactarlo de forma honesta y concreta.
